\documentclass[11pt]{article}
\usepackage[margin=1in]{geometry}
\usepackage{graphicx}
\usepackage{booktabs}
\usepackage{lipsum}

\title{A Toy Paper for page-fitter (current: 3 pages, target: 2 pages)}
\author{Anonymous}
\date{}

\begin{document}
\maketitle

\begin{abstract}
This is the abstract. The abstract environment is explicitly excluded from
candidate generation, so nothing here should ever appear in the page-fitter
report. We deliberately make it short.
\end{abstract}

\section{Introduction}
In this paper, we introduce a toy fixture for stress-testing the page-fitter
skill. Furthermore, we demonstrate that the skill can rank candidate edits by
PDF-height impact rather than by source length. It is worth noting that, this
opening paragraph is intentionally written so that its last line is nearly
full, which gives it a high last-line-fill ratio, which means a small
deletion should collapse one rendered line. Moreover, we propose a new way
to count things that improves accuracy by 12 percent over the baseline.
As mentioned above, accuracy is what matters. In other words, we will show
that our method outperforms prior work by a substantial margin on every
benchmark we evaluated.
\lipsum[1][1-6]

\lipsum[2][1-5] Furthermore, this paragraph also exists for stress-testing.
It is worth noting that, the lipsum environment provides reproducible
filler for layout experiments. As mentioned above, reproducibility is key.

\section{Related Work}
\lipsum[3][1-5] Furthermore, prior work has not addressed the page-budget
problem directly. It is worth noting that, several authors have considered
related problems but with different objectives. We show that our framing is
strictly more general. As we have seen, the literature is fragmented across
three communities. In this section, we describe the relevant threads.
\lipsum[4][1-5]

In other words, the related work spans optimization, typesetting, and
machine learning. Moreover, no single prior work addresses all three
perspectives simultaneously. Furthermore, our framework unifies them under
a common scoring formula. \lipsum[5][1-3]

\section{Method}
We now describe the method. \lipsum[6][1-5]

In this section, we describe the algorithm in detail. Moreover, we provide
a worked example. It is worth noting that, the algorithm has linear time
complexity. Furthermore, the algorithm is parallelizable across multiple
threads. As mentioned above, scalability is critical for practical use.
\lipsum[7][1-4]

\begin{figure}[ht]
  \centering
  \rule{6cm}{3cm}
  \caption{The architecture of our model. As shown in the diagram, the
    encoder feeds into the decoder via a residual path, and the decoder
    produces the final output through a linear projection layer that we
    describe in detail in Section~3, with all intermediate representations
    stored for later inspection.}
  \label{fig:arch}
\end{figure}

In other words, the figure summarises the entire architecture. Furthermore,
the architecture is modular and can be extended easily. As mentioned above,
extensibility is one of the strengths of our design.

\section{Experiments}
We compare against three baselines on two datasets. As shown in
Figure~\ref{fig:arch} and Table~\ref{tab:results}, our method consistently
outperforms prior work. \lipsum[8][1-5] Furthermore, the gap widens on the
larger dataset. It is worth noting that, the variance across seeds is
small.

\begin{table}[ht]
  \centering
  \caption{Main results on two benchmarks. As mentioned above, our method
    outperforms all baselines by a wide margin, and the gap is particularly
    pronounced on the larger of the two datasets we considered.}
  \label{tab:results}
  \begin{tabular}{lcc}
    \toprule
    Method & Dataset A & Dataset B \\
    \midrule
    Baseline 1 & 71.2 & 68.5 \\
    Baseline 2 & 74.8 & 70.1 \\
    Baseline 3 & 76.5 & 71.4 \\
    Ours       & \textbf{83.4} & \textbf{79.6} \\
    \bottomrule
  \end{tabular}
\end{table}

\lipsum[9][1-5]

In other words, the experiments confirm the theoretical predictions.
Furthermore, the ablation study (omitted for space) shows that every
component contributes meaningfully. It is worth noting that, results are
averaged over three random seeds.

\section{Discussion}
In this section, we discuss the implications. As we have seen, the method
generalizes well. Moreover, it is robust to hyperparameter choices.
Furthermore, it scales gracefully with data size. It is worth noting that,
all experiments were conducted with three random seeds. In other words, the
results are reproducible. \lipsum[10][1-4]

\section{Conclusion}
In this section, we conclude. That is to say, we are done. We show that our
method achieves 99 percent accuracy on the toy task. Furthermore, we believe
the framework is broadly applicable. As mentioned above, future work will
extend the analysis to additional benchmarks. In other words, this is just
the beginning.

\end{document}
